\section{From geometric to learned correspondence}
\label{sec:learned-correspondence}

Every correspondence mechanism considered so far in this report is
geometric: a vertex is matched to whichever target point is nearest to it,
possibly after the chain-aware and shape-diameter refinements of the previous
sections. A separate line of work asked whether a \emph{learned}
descriptor, a small network trained to produce a feature vector per
point such that corresponding anatomical locations produce similar vectors
could do better, and pursued that question through four increasingly
demanding tests.

The first attempt, a handcrafted baseline rather than a learned one,
normalised each point's coordinates within its own anatomical part (by that
part's local length and radius) before matching, on the hypothesis that
canonicalising away each part's specific proportions would make
correspondence easier to learn. It performed \emph{worse} than raw
coordinates: normalising away a part's specific geometry destroys
information that, it turns out, itself carries correspondence signal. The
first genuinely learned descriptor, trained with an off-the-shelf objective,
then failed when deployed, for a mundane but instructive reason: a mismatch
between the density of points it was trained on and the density it was
evaluated and deployed on, not a failure of the learning approach itself.
This motivated a redesigned, within-part training objective, tested next.

Two tasks must be kept apart. \emph{Pairwise retrieval} matches each query
point to its nearest neighbour in descriptor space,
$q_i\mapsto\arg\min_j\|f(q_i)-f(t_j)\|$. \emph{Dense correspondence} needs a
whole assignment $\{q_i\}_{i=1}^N\mapsto\{t_{\pi(i)}\}_{i=1}^N$ that is
globally consistent and covers the surface. Good retrieval does not
guarantee a consistent assignment.

\begin{algbox}{3: Learned-correspondence evaluation}
\begin{enumerate}[nosep,leftmargin=1.4em]
\item Compute descriptors $f$ for template and target points.
\item Retrieve candidates by nearest neighbour in descriptor space.
\item Score pairwise accuracy at a joint (retrieval).
\item Deploy the votes as a dense field over the whole mesh; measure coverage
and error against the geometric baseline across all specimens.
\end{enumerate}
\end{algbox}

\begin{investigation}{I14}{Can a descriptor trained specifically within one
anatomical part outperform raw local geometry at correspondence within that
part?}
\invQuestion{Trained and evaluated with an identical backbone and warm-start
procedure to the earlier attempt, does a within-part learned descriptor
improve correspondence accuracy at a specific, hard joint?}
\invMethod{Held-out coxa correspondence error, within-part learned
descriptor vs.\ the earlier learned baseline, same backbone, bootstrap
resampling for a confidence interval.}
\invResult{Coxa correspondence error falls from $0.2689$ to $0.2247$, a
$16.4\%$ improvement (bootstrap 95\% CI $[12.6\%, 19.8\%]$,
$p=7.8\times10^{-26}$), and the same descriptor outperforms the earlier
baseline across other segments too.}
\invVerdict{HOLD}{A learned local descriptor extracts correspondence
information that raw local geometry alone does not provide. This is
retained as evidence that learning has something to contribute, not as a
deployed mechanism: \S\ref{sec:learned-correspondence} next tests exactly
this descriptor as a dense correspondence field and rejects it at that
level. This is the strongest positive evidence in the investigation that
learning captures signal here, and, as the next result shows, it is also
where the harder problem begins.}
\end{investigation}

\begin{figure}[!htb]
\centering
\begin{tikzpicture}[every node/.style={font=\scriptsize},
 n/.style={draw,rounded corners=2pt,minimum height=7mm,align=center,fill=white,inner sep=3pt}]
\node[n,fill=green!12] (a) {Local test\\checkmark};
\node[n,below=3mm of a] (a2) {within-part descriptor};
\node[n,below=3mm of a2] (a3) {better pairwise retrieval};
\node[n,below=3mm of a3,fill=green!12] (a4) {$+16.4\%$ coxa};
\node[n,right=22mm of a,fill=red!10] (b) {Deployment\\$\times$};
\node[n,below=3mm of b] (b2) {dense correspondence};
\node[n,below=3mm of b2] (b3) {votes collapse};
\node[n,below=3mm of b3,fill=red!10] (b4) {$28.1\%\to1.8\%$ coverage};
\node[below=4mm of $(a4)!0.5!(b4)$,font=\small\bfseries] {local success $\neq$ global consistency};
\end{tikzpicture}
\caption{The same descriptor, evaluated at two levels of abstraction. Strong
local (pairwise) retrieval does not imply a usable dense correspondence
field.}
\label{fig:local-vs-global}
\end{figure}

\begin{investigation}{I15}{Does strong pairwise retrieval accuracy translate
into a usable, dense correspondence field at deployment?}
\invQuestion{Deployed as a dense, whole-mesh correspondence mechanism rather
than evaluated pairwise at one joint, does the within-part descriptor still
outperform the existing geometric baseline?}
\invMethod{Full-mesh deployment benchmark across 48 specimens, learned
descriptor vs.\ geometric nearest-point baseline; a follow-up analysis
isolating how much of any gap is attributable to differences in the
candidate point set each method was allowed to match against, versus the
descriptor's votes themselves.}
\invResult{The learned descriptor performs \emph{worse} at deployment scale.
Candidate-set differences account for only about $2.9\%$ of the observed
gap, against an overall gap of roughly $145\%$; at the coxa specifically,
surface coverage collapses from $28.1\%$ under the geometric baseline to
$1.8\%$ under the learned descriptor.}
\invVerdict{REJECT}{The learned votes collapse onto a small subset of
points instead of producing a dense, globally consistent field across the
whole mesh. Pairwise retrieval accuracy at a single joint (Investigation~I14) is not
the same property as dense, globally consistent correspondence across an
entire specimen, and the two must be evaluated separately rather than
assumed to track each other.}
\end{investigation}

Read together, these four steps identify a specific, narrow gap rather than
a blanket failure of learning: local, within-part descriptors carry real
information (Investigation~I14) that geometric nearest-matching does not,
but turning that local signal into a dense, mesh-wide correspondence
field (global consistency, not just local accuracy) is an unsolved
engineering problem in its own right (Investigation~I15), and is the
concrete direction this report's outlook (Chapter~\ref{ch:conclusion})
recommends pursuing further.
